\documentclass[11pt,a4paper]{article}

% ── Codificación y fuentes ────────────────────────────────────────────────────
\usepackage[utf8]{inputenc}
\usepackage[T1]{fontenc}
\usepackage{lmodern}
\usepackage[spanish]{babel}

% ── Geometría y espaciado ─────────────────────────────────────────────────────
\usepackage[top=2.5cm, bottom=2.5cm, left=2.8cm, right=2.8cm]{geometry}
\usepackage{setspace}
\setstretch{1.15}
\usepackage{parskip}

% ── Tablas ────────────────────────────────────────────────────────────────────
\usepackage{booktabs}
\usepackage{tabularx}
\usepackage{array}
\usepackage{longtable}
\usepackage{enumitem}

% ── Colores y cajas ───────────────────────────────────────────────────────────
\usepackage[dvipsnames]{xcolor}
\usepackage{tcolorbox}
\tcbuselibrary{skins, breakable}

% ── Cabeceras y pies de página ────────────────────────────────────────────────
\usepackage{fancyhdr}
\usepackage{lastpage}
\pagestyle{fancy}
\fancyhf{}
\fancyhead[L]{\small\textbf{Informe de Evaluación} --- Física para Bioanalistas}
\fancyhead[R]{\small valeria marrero 34264302}
\fancyfoot[C]{\small Página \thepage\ de \pageref{LastPage}}
\renewcommand{\headrulewidth}{0.4pt}

% ── Hiperenlaces ──────────────────────────────────────────────────────────────
\usepackage[colorlinks=true, linkcolor=NavyBlue, urlcolor=NavyBlue]{hyperref}

% ── Paleta de colores personalizados ─────────────────────────────────────────
\definecolor{desempeno_excelente}{HTML}{1A6B2F}
\definecolor{desempeno_bueno}{HTML}{2E5A9C}
\definecolor{desempeno_suficiente}{HTML}{8B6914}
\definecolor{desempeno_deficiente}{HTML}{C0392B}
\definecolor{desempeno_insuficiente}{HTML}{7B241C}
\definecolor{grisFondo}{HTML}{F5F5F5}
\definecolor{azulTitulo}{HTML}{1B2A6B}
\definecolor{verdeFortaleza}{HTML}{1A6B2F}
\definecolor{naranjaMejora}{HTML}{A04000}

% ── Estilos de cajas ─────────────────────────────────────────────────────────
\tcbset{
  cajaTitulo/.style={
    enhanced, breakable,
    colback=azulTitulo!8, colframe=azulTitulo,
    fonttitle=\bfseries\large, coltitle=white,
    attach boxed title to top left={yshift=-2mm, xshift=4mm},
    boxed title style={colback=azulTitulo},
    top=6pt, bottom=6pt, left=8pt, right=8pt,
  },
  cajaFortaleza/.style={
    enhanced, breakable,
    colback=verdeFortaleza!6, colframe=verdeFortaleza!60,
    fonttitle=\bfseries, coltitle=verdeFortaleza,
    top=4pt, bottom=4pt, left=8pt, right=8pt,
  },
  cajaMejora/.style={
    enhanced, breakable,
    colback=naranjaMejora!6, colframe=naranjaMejora!60,
    fonttitle=\bfseries, coltitle=naranjaMejora,
    top=4pt, bottom=4pt, left=8pt, right=8pt,
  },
  cajaRecomendacion/.style={
    enhanced, breakable,
    colback=grisFondo, colframe=gray!50,
    fonttitle=\bfseries, coltitle=black,
    top=4pt, bottom=4pt, left=8pt, right=8pt,
  },
}

% ─────────────────────────────────────────────────────────────────────────────
\begin{document}

% ══ PORTADA ══════════════════════════════════════════════════════════════════
\begin{center}
  {\Large\bfseries\color{azulTitulo} Universidad de Oriente/Núcleo Sucre --- Física para Bioanalistas}\\[4pt]
  {\large Informe de Evaluación Preliminar}\\[12pt]
  \rule{\linewidth}{1.2pt}\\[6pt]
  {\LARGE\bfseries valeria marrero 34264302}\\[4pt]
  \rule{\linewidth}{0.6pt}
\end{center}

% ══ FICHA TÉCNICA ════════════════════════════════════════════════════════════
\vspace{4pt}
\begin{tcolorbox}[cajaTitulo, title=Datos del informe]
\begin{tabular}{@{}llll@{}}
  \textbf{Estudiante:}  & valeria marrero 34264302  &
  \textbf{Fecha:}       & 2026-06-26 \\[4pt]
  \textbf{Archivo PDF:} & \texttt{valeria\_marrero\_34264302.pdf}  &
  \textbf{Páginas:}     & 4 \\
\end{tabular}
\end{tcolorbox}

% ══ RESULTADO GLOBAL ═════════════════════════════════════════════════════════
\vspace{6pt}
\begin{center}
  \begin{tcolorbox}[
    enhanced,
    colback=white, colframe=desempeno_bueno,
    width=0.62\linewidth,
    boxrule=2pt,
    halign=center, valign=center,
    top=8pt, bottom=8pt,
  ]
    \centering
    {\large\bfseries Puntaje sugerido}\\[6pt]
    {\Huge\bfseries\color{desempeno_bueno} 7.3/10}\\[6pt]
    {\large Nivel de desempeño: \textbf{\color{desempeno_bueno} Bueno}}
  \end{tcolorbox}
\end{center}

% ══ RESUMEN GENERAL ══════════════════════════════════════════════════════════
\section*{Resumen general}
El estudiante demuestra un buen entendimiento de los principios físicos fundamentales de la hidrostática, la hemodinámica y la ley de Poiseuille. Es capaz de identificar las fórmulas correctas y organizar sus soluciones de manera clara. Sin embargo, se observan errores recurrentes en la precisión de los cálculos matemáticos, especialmente en el manejo de exponentes y en la correcta aplicación de cantidades derivadas en las fórmulas, lo que afecta la exactitud de los resultados finales en algunos problemas.

% ══ FORTALEZAS Y ÁREAS DE MEJORA ════════════════════════════════════════════
\vspace{4pt}
\begin{minipage}[t]{0.48\linewidth}
  \begin{tcolorbox}[cajaFortaleza, title={Fortalezas}]
\begin{itemize}[leftmargin=1.5em, itemsep=2pt]
  \item Comprensión de los principios físicos subyacentes en la mayoría de los problemas.
  \item Habilidad para extraer datos relevantes y realizar conversiones de unidades correctamente.
  \item Claridad y organización en la presentación de los pasos de la solución.
  \item Correcta aplicación de la ecuación de continuidad y la ley de Poiseuille en su forma proporcional.
\end{itemize}
  \end{tcolorbox}
\end{minipage}
\hfill
\begin{minipage}[t]{0.48\linewidth}
  \begin{tcolorbox}[cajaMejora, title={Areas de mejora}]
\begin{itemize}[leftmargin=1.5em, itemsep=2pt]
  \item Precisión en los cálculos numéricos, incluyendo el manejo de exponentes.
  \item Aplicación rigurosa de las definiciones y fórmulas, evitando confusiones entre magnitudes relacionadas (ej. fuerza de empuje vs. peso aparente).
  \item Revisión de los pasos intermedios para detectar y corregir errores de cálculo.
\end{itemize}
  \end{tcolorbox}
\end{minipage}

% ══ OBSERVACIONES POR PREGUNTA ═══════════════════════════════════════════════
\section*{Observaciones por pregunta}

\begin{longtable}{>{\bfseries}p{2.8cm} p{1.6cm} p{7.0cm} p{2.8cm}}
  \toprule
  \textbf{Pregunta} & \textbf{Puntaje} & \textbf{Evaluación} & \textbf{Errores detectados} \\
  \midrule
  \endhead
  \bottomrule
  \endfoot
    \textbf{1. Presión Hidrostática y Venosa} & 2.0/2.0 & El estudiante interpretó correctamente la pregunta como la búsqueda de la altura 'h' necesaria para generar la presión manométrica dada. La fórmula utilizada es la adecuada y el cálculo es correcto, incluyendo la conversión de unidades y el resultado final. & \textit{---} \\
    \midrule
    \textbf{2. Principio de Arquímedes} & 0.5/2.0 & El estudiante calculó correctamente la fuerza de empuje (E) y la masa del hueso. Sin embargo, cometió un error conceptual al utilizar el peso aparente (2.40 N) en lugar de la fuerza de empuje (2.10 N) para calcular el volumen del objeto sumergido. Esto llevó a un volumen incorrecto y, consecuentemente, a una densidad incorrecta. Además, hubo un error de cálculo en la densidad final (4873.47 kg/m$^{3}$ no corresponde a 0.459 / 2.45x10$^{-}$$^{4}$). & \textit{Error conceptual: Uso incorrecto del peso aparente en lugar de la fuerza de empuje para calcular el volumen.; Error procedimental: Cálculo incorrecto de la densidad final.} \\
    \midrule
    \textbf{3. Hemodinámica (Continuidad y Bernoulli)} & 1.8/2.0 & El estudiante aplicó correctamente la ecuación de continuidad para calcular la velocidad en la zona de estenosis (v2). La aplicación de la ecuación de Bernoulli es correcta, asumiendo un flujo horizontal. Se detectó un error de cálculo menor en un paso intermedio (530 * 1.35 = 715.5, pero el estudiante escribió 775.5). A pesar de este error intermedio, el resultado final para P2 (13284.5 Pa) es correcto, lo que sugiere una posible corrección posterior o un error de transcripción en el paso intermedio. & \textit{Error procedimental: Cálculo incorrecto en un paso intermedio (530 * 1.35).} \\
    \midrule
    \textbf{4. Ley de Poiseuille (Análisis Proporcional)} & 2.0/2.0 & El estudiante identificó correctamente la relación de proporcionalidad entre el caudal (Q) y el radio a la cuarta potencia (r$^{4}$) según la Ley de Poiseuille. El cálculo de la reducción del radio y su impacto en el caudal es preciso. La interpretación de la pregunta sobre el 'porcentaje de caudal que se mantendrá circulando' es adecuada, y el resultado del 40.96\% es correcto. & \textit{---} \\
    \midrule
    \textbf{5. Flujo Viscoso y Resistencia} & 1.0/2.0 & El estudiante identificó correctamente la Ley de Poiseuille y la despejó para calcular la diferencia de presión ($\Delta$P). La extracción de datos y la conversión de unidades son correctas. Sin embargo, se cometió un error significativo en el cálculo de r$^{4}$: (4 x 10$^{-}$$^{4}$ m)$^{4}$ es 2.56 x 10$^{-}$$^{1}$$^{4}$ m$^{4}$, no 2.56 x 10$^{-}$$^{1}$$^{3}$ m$^{4}$. Este error en el exponente propagó un error de un factor de 10 en el resultado final de $\Delta$P. & \textit{Error procedimental: Cálculo incorrecto del exponente de r$^{4}$, lo que llevó a un resultado final erróneo por un factor de 10.} \\
    \midrule
\end{longtable}

% ══ ERRORES DETECTADOS ═══════════════════════════════════════════════════════
\section*{Análisis de errores}

\subsection*{Errores conceptuales}
\begin{itemize}[leftmargin=1.5em, itemsep=2pt]
  \item Confusión entre la fuerza de empuje y el peso aparente al calcular el volumen de un objeto sumergido (Pregunta 2).
\end{itemize}

\subsection*{Errores procedimentales}
\begin{itemize}[leftmargin=1.5em, itemsep=2pt]
  \item Errores de cálculo en la determinación de la densidad (Pregunta 2).
  \item Error menor de cálculo en un paso intermedio (Pregunta 3).
  \item Error significativo en el cálculo de un exponente (r$^{4}$) que afectó el resultado final por un factor de 10 (Pregunta 5).
\end{itemize}

% ══ RECOMENDACIONES AL ESTUDIANTE ════════════════════════════════════════════
\section*{Retroalimentación para el estudiante}
\begin{tcolorbox}[cajaRecomendacion, title={Dirigido a: valeria marrero 34264302}]
Estimado/a estudiante, tu examen muestra un buen dominio de los conceptos físicos en general. Has demostrado la capacidad de identificar las fórmulas correctas y organizar tus soluciones de manera clara. Sin embargo, es crucial que prestes más atención a la precisión en tus cálculos, especialmente al manejar exponentes y al sustituir valores en las fórmulas. Por ejemplo, en el problema de Arquímedes, asegúrate de distinguir entre la fuerza de empuje y el peso aparente al calcular el volumen. En el problema de Poiseuille, un pequeño error en un exponente afectó significativamente el resultado final. Te sugiero revisar tus operaciones matemáticas con mayor detalle y verificar tus resultados. ¡Sigue practicando para fortalecer estas habilidades!
\end{tcolorbox}

% ══ NOTA PARA EL PROFESOR ════════════════════════════════════════════════════
\section*{Nota para el profesor}
\begin{tcolorbox}[
  enhanced, breakable,
  colback=yellow!8, colframe=yellow!60!black,
  fonttitle=\bfseries, coltitle=black,
  title={Observaciones para revisión manual},
  top=4pt, bottom=4pt, left=8pt, right=8pt,
]
El estudiante demuestra una comprensión sólida de los principios de hidrostática, continuidad y la Ley de Poiseuille. Los errores principales se encuentran en la ejecución de cálculos y la aplicación precisa de las cantidades derivadas (como la fuerza de empuje en el problema 2). En el problema 3, hay un error numérico en un paso intermedio (530 * 1.35 = 775.5 en lugar de 715.5), pero el resultado final es curiosamente correcto (13284.5 Pa), lo que podría indicar una corrección posterior o un error de transcripción. El problema 5 tiene un error significativo en el cálculo de r$^{4}$ (exponente incorrecto), lo que lleva a un resultado final erróneo por un factor de 10. Se recomienda enfatizar la importancia de la precisión matemática y la revisión de los pasos intermedios.
\end{tcolorbox}

% ══ PIE DE INFORME ════════════════════════════════════════════════════════════
\vspace{12pt}
\begin{center}
  \footnotesize\color{gray}
  Informe generado automáticamente por el sistema de evaluación con IA (gemini-2.5-flash) y soporte LaTex. \\
  Este documento es una evaluación \textbf{preliminar} para ser revisado por el profesor antes de ser entregado al 
  estudiante. \\
  Generado el 2026-06-26 \textbullet\ BIOANALISIS Sistema Académico de Evaluación.
  \end{center}

\end{document}
